\section{Conclusion}
We presented \framework, an end-to-end evaluation framework for voice agents that jointly addresses simulation fidelity and measurement comprehensiveness. Together, validation-gated bot-to-bot simulation, architecture-agnostic composite metrics (EVA-A and EVA-X), and a multi-trial consistency framework (\passatone, \passatk, \passpowerk) enable comparison across cascade, hybrid, and S2S voice agents under identical conditions.

Our evaluation of 12 systems produces three central findings. First, while the best-performing cascade and S2S systems achieve comparable accuracy, experience quality diverges sharply along architecture lines,  with the S2S--cascade gap on EVA-X driven almost entirely by turn-taking. Second, peak and reliable performance diverge substantially across all systems: the median \passatk--\passpowerk~gap is 0.44 on EVA-A, indicating that single-trial evaluation scores systematically overstate deployment-grade quality regardless of architecture. Third, we observe that cascade and S2S systems degrade asymmetrically under acoustic perturbation: accent variation degrades the cascade task completion by an average of 10 points while leaving S2S accuracy unchanged, whereas background noise degrades experience metrics for almost all systems. These observations highlight dimensions of voice agent quality that existing benchmarks leave unmeasured: joint accuracy-experience measurement, peak-versus-reliable performance gaps, and perturbation-specific degradation patterns across voice agent architectures. 

\framework~is released as a fully open-source, extensible framework. The metric 
suite is modular, the simulation interface supports cascade and audio-native 
architectures, and new domains can be added by configuring agents and scenarios. We invite the community to extend the benchmark with additional domains, languages, and evaluation dimensions as the field evolves.
